% ============================================================================
% Fichier  : partie3/P03_C13_conduire.tex
% Titre    : Conduire l'Audit : Preuves, Échantillons, Constats
% Auteur   : MINKA MI NGUIDJOÏ Thierry Emmanuel
% Version  : 1.0 -- novembre 2026
% ============================================================================
% Sources : ISO 19011:2018 (6.3 à 6.4, annexe A) ; ISO/IEC 27007:2020.
% Tailles d'échantillon et limites supérieures : loi binomiale, confiance
% 95 %, calculées et vérifiées. Audit d'octobre 2026 : cas fictif.
% ============================================================================

\chapter{Conduire l'Audit : Preuves, Échantillons, Constats}
\label{chap:conduire}
\chaptermark{Preuves, échantillons, constats}

\epigraph{%
    « Zéro anomalie dans un échantillon\\
    ne veut pas dire zéro anomalie dans la population. »%
}{Principe du cours}

\niveauChapitre{Intermédiaire à expert}{Chapitre~\ref{chap:audit}~;
notions de probabilité.}

\begin{boxobjectifs}
\begin{itemize}
    \item Dérouler les étapes d'un audit, de la préparation à la réunion de
          clôture.
    \item Juger la force d'une preuve selon sa source.
    \item Dimensionner un échantillon par attributs et en tirer une limite
          supérieure d'anomalies.
    \item Rédiger un constat~: critère, fait, preuve, cause, conséquence.
    \item Distinguer non-conformité, opportunité d'amélioration et
          conformité.
\end{itemize}
\end{boxobjectifs}

\bigskip

% ============================================================================
\section*{Question directrice : 29 connexions examinées suffisent-elles à conclure sur 1 340~?}
% ============================================================================

En octobre, Mme~NKENG audite les accès distants de LiZbiZ. En septembre,
la passerelle a enregistré 1\,340~connexions. Elle ne peut pas les examiner
toutes à la main. Combien en examiner, comment les choisir, et que
conclure si elle trouve une ou deux connexions sans second facteur~?

\medskip
\textbf{Plan du chapitre.}
\begin{enumerate}
    \item Les étapes d'un audit (section~\ref{sec:c13-etapes}).
    \item Les preuves (section~\ref{sec:c13-preuves}).
    \item Échantillonner (section~\ref{sec:c13-echantillon}).
    \item Conclure sur la population (section~\ref{sec:c13-conclure}).
    \item Rédiger un constat (section~\ref{sec:c13-constat}).
\end{enumerate}

\separateur

% ============================================================================
\section{Les Étapes d'un Audit}
\label{sec:c13-etapes}
% ============================================================================

\begin{tableaucours}{L{3.6cm}L{8.6cm}}{%
    \tableauHeader{Étape} &
    \tableauHeader{Contenu (d'après l'ISO 19011:2018)}
}
Déclenchement & Contact avec l'audité, faisabilité, accès aux informations \\
Préparation & Revue documentaire, plan d'audit, répartition des tâches,
documents de travail (listes de vérification, plans d'échantillonnage) \\
Réunion d'ouverture & Confirmer le périmètre, les critères, le
calendrier, les modalités \\
Collecte et vérification & Entretiens, observations, examen de documents et
d'enregistrements, essais \\
Constats & Comparer les preuves aux critères \\
Conclusions & Synthèse des constats au regard des objectifs de l'audit \\
Réunion de clôture & Présenter constats et conclusions, recueillir les
désaccords \\
Rapport et suivi & Rapport, actions correctives, vérification
(Chapitre~\ref{chap:rapport}) \\
\end{tableaucours}
\vspace{-0.8em}
\begin{center}{\small\itshape Tableau~13.1. Les étapes d'un audit}\end{center}

% ============================================================================
\section{Les Preuves}
\label{sec:c13-preuves}
% ============================================================================

\begin{boxdef}
\termecle{Preuve d'audit}~: enregistrements, énoncés de faits ou autres
informations, pertinents pour les critères d'audit et vérifiables.
\end{boxdef}

\begin{tableaucours}{L{3.4cm}L{2.2cm}L{6.6cm}}{%
    \tableauHeader{Source} &
    \tableauHeader{Force} &
    \tableauHeader{Exemple}
}
Essai refait par l'auditeur & Élevée & Tenter une connexion à la passerelle
sans second facteur \\
Enregistrement produit par le système & Élevée & Journal de la passerelle,
exporté en présence de l'auditeur, avec son empreinte \\
Document & Moyenne & Politique d'accès distant signée \\
Observation & Moyenne & Voir le technicien se connecter par le bastion \\
Déclaration & Faible & « Tout le monde a le second facteur depuis juillet » \\
\end{tableaucours}
\vspace{-0.8em}
\begin{center}{\small\itshape Tableau~13.2. Force des preuves selon leur source}\end{center}

\begin{boxalert}
\textbf{Violation typique.} Conclure sur une déclaration.

\smallskip
\textbf{Règle~:} une déclaration oriente la recherche, elle ne prouve pas.
On la corrobore par un enregistrement ou un essai. Faute de corroboration,
le point reste ouvert (au Chapitre~\ref{chap:retex}, on le cotera \NEK{}).
\end{boxalert}

% ============================================================================
\section{Échantillonner}
\label{sec:c13-echantillon}
% ============================================================================

Quand la population est trop grande, l'auditeur examine un
\termecle{échantillon}. Il peut le choisir \emph{par jugement} (les
connexions de nuit, celles du prestataire) ou \emph{statistiquement}
(au hasard), ce qui seul permet d'étendre une conclusion chiffrée à la
population.

\begin{boxdef}
\termecle{Échantillonnage par attributs}~: on compte, dans un échantillon
tiré au hasard, les éléments qui présentent une anomalie (oui ou non). On
fixe un \emph{taux tolérable} d'anomalies et un \emph{niveau de confiance}.
\end{boxdef}

Si l'on s'attend à ne trouver aucune anomalie, la taille minimale $n$ pour
conclure, avec une confiance de 95~\%, que le taux d'anomalies est
inférieur au taux tolérable $t$ vérifie $(1 - t)^n \le 0{,}05$~:
\[
n \ge \frac{\ln 0{,}05}{\ln(1 - t)}.
\]

\begin{tableaucours}{L{4.0cm}R{3.0cm}}{%
    \tableauHeader{Taux tolérable} &
    \tableauHeader{Taille minimale}
}
10~\% & 29 \\
5~\% & 59 \\
3~\% & 99 \\
\end{tableaucours}
\vspace{-0.8em}
\begin{center}{\small\itshape Tableau~13.3. Tailles d'échantillon, confiance 95~\%, aucune anomalie attendue}\end{center}

La taille dépend du taux tolérable et de la confiance, presque pas de la
taille de la population, dès que celle-ci est grande. 29 connexions
suffisent pour 1\,340 comme pour 100\,000.

% ============================================================================
\section{Conclure sur la Population}
\label{sec:c13-conclure}
% ============================================================================

Si l'échantillon contient des anomalies, on calcule la \termecle{limite
supérieure} du taux d'anomalies dans la population, à 95~\% de confiance.

\begin{tableaucours}{L{3.6cm}R{2.8cm}R{2.8cm}}{%
    \tableauHeader{Anomalies trouvées} &
    \tableauHeader{Échantillon de 29} &
    \tableauHeader{Échantillon de 59}
}
0 & 9,8~\% & 5,0~\% \\
1 & 15,3~\% & 7,8~\% \\
2 & 20,2~\% & 10,3~\% \\
\end{tableaucours}
\vspace{-0.8em}
\begin{center}{\small\itshape Tableau~13.4. Limite supérieure du taux d'anomalies, confiance 95~\%}\end{center}

\begin{boxscenario}{Opération LIZBIZ. Les connexions de septembre}
\textbf{Constatation.} Sur 29~connexions tirées au hasard parmi 1\,340,
une s'est faite sans second facteur, par le compte « vpn\_secours »,
exclu de la politique pour « dépannage ». L'essai de l'auditrice confirme
que ce compte se connecte avec son seul mot de passe.

\medskip
\textbf{Analyse} (confiance élevée). Avec une anomalie sur 29, le taux
réel peut atteindre 15,3~\%~: on ne peut pas conclure que 90~\% au moins
des connexions ont un second facteur. Mais l'enjeu n'est plus statistique~:
la cause est identifiée (un compte exclu), et toutes les connexions de ce
compte sont à examiner, par un test exhaustif ciblé.
\end{boxscenario}

% ============================================================================
\section{Rédiger un Constat}
\label{sec:c13-constat}
% ============================================================================

\begin{boxprocedure}
\textbf{Les éléments d'un constat}
\begin{enumerate}
    \item \textbf{Critère}~: l'exigence (politique, clause, article).
    \item \textbf{Constat}~: le fait observé, chiffré.
    \item \textbf{Preuve}~: la pièce qui l'établit, avec sa référence.
    \item \textbf{Cause}~: pourquoi l'écart existe.
    \item \textbf{Conséquence}~: le risque qui en résulte, relié au registre.
    \item \textbf{Classement}~: non-conformité (majeure ou mineure),
          opportunité d'amélioration, ou conformité.
\end{enumerate}
\end{boxprocedure}

\begin{boxobjet}{Constat A-03 · Compte d'accès distant sans second facteur}
\textbf{Critère~:} politique d'accès distant (juillet 2026)~: « toute
connexion distante exige un second facteur ». \textbf{Constat~:} le compte
« vpn\_secours » se connecte sans second facteur~; 46~connexions en
septembre. \textbf{Preuve~:} journal de la passerelle (export du 8~octobre,
empreinte SHA-256 consignée), essai du 9~octobre. \textbf{Cause~:}
exception accordée oralement en juillet, jamais revue. \textbf{Conséquence~:}
R-02 n'est pas réduit comme le registre l'indique. \textbf{Classement~:}
non-conformité majeure.
\end{boxobjet}

\begin{boxmetier}{Auditeur des systèmes d'information}
\textbf{Chiffrer, puis expliquer.} L'auditeur ne se contente pas d'écrire
« certaines connexions »~: il donne le nombre, la période, la preuve, et
passe du sondage au test exhaustif dès qu'il a identifié une cause.
\end{boxmetier}

\separateur

\begin{boxresume}
\textbf{Ce qu'il faut retenir de ce chapitre~:}
\begin{itemize}
    \item \textbf{Un audit} se prépare, s'ouvre, collecte, constate,
          conclut et se clôt en réunion.
    \item \textbf{Une preuve} vaut selon sa source~: essai et enregistrement
          avant document, observation et déclaration.
    \item \textbf{L'échantillon statistique} seul permet d'étendre une
          conclusion~: 29 pour 10~\% de taux tolérable à 95~\%.
    \item \textbf{Une anomalie trouvée} élève la limite supérieure~; une
          cause identifiée appelle un test exhaustif.
    \item \textbf{Un constat}~: critère, fait, preuve, cause, conséquence,
          classement.
\end{itemize}
\end{boxresume}

\begin{boxinfo}
\textbf{Points de contrôle de l'élève.} Je sais~: dérouler les étapes d'un
audit~; classer une preuve par sa force~; calculer une taille d'échantillon
sans anomalie attendue~; lire une limite supérieure~; rédiger un constat
complet.
\end{boxinfo}

\bigskip

\begin{boxexo}[title={\ding{45}\quad Exercice~13.1. Les étapes \niveaubase}]
À quelle étape de l'audit se fait~: la confirmation du périmètre avec
l'audité~; l'export du journal~; la présentation des constats~?
\end{boxexo}

\begin{boxexo}[title={\ding{45}\quad Exercice~13.2. Classer les preuves \niveaubase}]
Classez par force décroissante~: un courriel du prestataire~; l'export
du journal~; la politique signée~; un essai de connexion.
\end{boxexo}

\begin{boxexo}[title={\ding{45}\quad Exercice~13.3. Taille \niveaubase}]
Quelle taille d'échantillon faut-il pour conclure, à 95~\% de confiance et
sans anomalie attendue, à un taux d'anomalies inférieur à 5~\%~?
\end{boxexo}

\begin{boxexo}[title={\ding{45}\quad Exercice~13.4. La formule \niveaumoyen}]
Vérifiez par le calcul que 29 est la plus petite taille qui satisfait
$(1 - 0{,}10)^n \le 0{,}05$.
\end{boxexo}

\begin{boxexo}[title={\ding{45}\quad Exercice~13.5. Lire la limite \niveaumoyen}]
Avec deux anomalies dans un échantillon de 59, peut-on conclure à un taux
inférieur à 10~\%~? Justifiez par le Tableau~13.4.
\end{boxexo}

\begin{boxexo}[title={\ding{45}\quad Exercice~13.6. Par jugement \niveaumoyen}]
Dans quels cas l'auditeur doit-il préférer un échantillon par jugement~?
Que ne peut-il alors pas conclure~?
\end{boxexo}

\begin{boxexo}[title={\ding{45}\quad Exercice~13.7. Le constat \niveaumoyen}]
Rédigez un constat complet pour l'écart suivant~: la sauvegarde de la
facturation n'a pas été testée depuis sa mise en place en août.
\end{boxexo}

\begin{boxexo}[title={\ding{45}\quad Exercice~13.8. Majeure ou mineure \niveaumoyen}]
Classez en non-conformité majeure, mineure ou opportunité
d'amélioration~: une politique non signée~; aucun registre des dérogations~;
le compte « vpn\_secours ».
\end{boxexo}

\begin{boxexo}[title={\ding{45}\quad Exercice~13.9. Le plan d'audit \niveauexpert}]
Rédigez le plan de l'audit « accès distants et prestataire »~: objectifs,
critères, périmètre, calendrier, échantillons, preuves attendues.
\end{boxexo}

\begin{boxexo}[title={\ding{45}\quad Exercice~13.10. Calculer une limite \niveauexpert}]
Écrivez la méthode (ou le programme) qui calcule la limite supérieure à
95~\% du taux d'anomalies pour $k$ anomalies dans un échantillon de $n$,
et vérifiez la valeur 15,3~\% pour $k=1$, $n=29$.
\end{boxexo}

\bigskip

\begin{boxinfo}
\textbf{Pour aller plus loin.}
\begin{itemize}
    \item ISO~19011, conduite de l'audit~\cite{iso19011_2018}.
    \item ISO/IEC~27007~\cite{iso27007_2020}.
\end{itemize}
\end{boxinfo}

\begin{boxinfo}
\textbf{Transition.} Le Chapitre~\ref{chap:auditer-risques} applique la
méthode à un objet particulier~: le dispositif de gestion des risques et
le système de management de la sécurité.
\end{boxinfo}
